\chapter{High-Frequency Event Loops}
\label{ch:high_frequency_event_loops}
\label{ch:event_loops}

Managing 400 autonomous Large Language Model agents as they asynchronously evaluate market states, compute semantic trajectories, and post limit orders presents a severe orchestration challenge. This chapter details the departure from standard Python multithreading in favor of cooperative multitasking via \texttt{asyncio}, the integration of the \texttt{uvloop} C-extension, and the implementation of Singleflight request coalescing to prevent catastrophic GPU cache stampedes.

\section{The Multithreading Context-Switch Penalty}

A naive implementation of AMPS might spawn a unique operating system thread for every generative agent. With 400 agents, the Linux kernel must continuously pause and resume these 400 threads to give them all time on the physical CPU cores. 

Every time the kernel switches from one thread to another, it incurs a \textit{Context-Switch Penalty}. The CPU must save the exact state of its registers, clear the pipeline, load the new thread's registers, and warm up the L1 cache. With 400 threads constantly demanding execution, the kernel wastes enormous amounts of CPU cycles merely juggling the threads, utterly destroying the simulation's throughput and violating our microsecond latency budgets.

\section{Cooperative Multitasking via \texttt{asyncio} and \texttt{uvloop}}

To resolve this, AMPS abandons OS-level multithreading entirely in favor of \textbf{Cooperative Multitasking}. Python's \texttt{asyncio} module operates on a single OS thread. Instead of the kernel forcibly interrupting agents, agents explicitly yield control of the CPU when they are waiting for I/O (such as waiting for the GPU to finish LLM inference or waiting for a timer to expire).

While the default Python \texttt{asyncio} event loop is a massive improvement over multithreading, it is written in pure Python and remains relatively slow. AMPS replaces the default loop with \textbf{\texttt{uvloop}} \cite{selivanov2016uvloop}.

\texttt{uvloop} is a Cython wrapper around \texttt{libuv} (the ultra-fast C library that powers Node.js). Under the hood, \texttt{uvloop} hooks directly into the operating system's lowest-level I/O notification facilities -- specifically the \texttt{epoll} system call on Linux. This paradigm shift allows the single Python thread to monitor thousands of agent states simultaneously in the C-kernel with practically zero polling overhead.

\section{Thundering Herd Mitigation (Singleflight Coalescing)}

Even with \texttt{uvloop}, AMPS faces a critical systems engineering crisis during a simulated market panic. 

When the 671B MoE Macro God Agent broadcasts a severe, systemic market shock (e.g., ``Major Bank Files for Bankruptcy''), all 400 micro-agents wake up simultaneously. They all instantly attempt to query the local RTX 5070 Ti GPU for generative inference to decide whether to panic sell. This instantaneous, massive spike in I/O demand is known as a \textbf{Cache Stampede} or a "Thundering Herd."

Sending 400 identical LLM prompts to the GPU simultaneously will exhaust the 16GB of VRAM and crash the CUDA kernels. To mitigate this fatal bottleneck, AMPS implements the \textbf{Singleflight} request coalescing pattern \cite{colakel2026singleflight}.

\begin{figure}[htbp]
\centering
\begin{tikzpicture}[
    scale=0.9, transform shape,
    node distance=1.5cm,
    agent/.style={circle, draw, fill=blue!10, minimum size=0.8cm},
    lock/.style={rectangle, draw, fill=red!10, rounded corners, minimum height=1cm, align=center},
    gpu/.style={rectangle, draw, fill=green!10, minimum height=1cm, align=center}
]
    \node[agent] (a1) {A1};
    \node[agent] (a2) [below=0.2cm of a1] {A2};
    \node[agent] (a3) [below=0.2cm of a2] {A3};
    
    \node[lock] (mutex) [right=1.5cm of a2] {Singleflight\\Mutex Map};
    \node[gpu] (gpu) [right=1.5cm of mutex] {RTX 5070 Ti\\(Inference)};
    
    \draw[->, thick, >=stealth] (a1) -- node[above=5pt, scale=0.8] {1st Caller} (mutex);
    \draw[->, dashed] (a2) -- node[below, scale=0.8] {Waits} (mutex);
    \draw[->, dashed] (a3) -- node[below=4pt, scale=0.8] {Waits} (mutex);
    
    \draw[->, thick, >=stealth] (mutex) -- node[above, scale=0.8] {Compute} (gpu);
    
    \draw[->, thick, >=stealth, blue] (gpu.south) to[bend left=30] node[below=1.5pt, scale=0.8,pos=0.55] {$O(1)$ Broadcast Promise} (a3.south east);
\end{tikzpicture}
\caption{Mechanics of Singleflight coalescing. The first agent locks the operation and executes it. All subsequent agents yield their context and await the broadcast result, completely flattening GPU utilization spikes.}
\label{fig:singleflight_detail}
\end{figure}

\section{Two-Phase Inference Pipeline: Resolving Agent State Heterogeneity}

A fundamental architectural pitfall in applying request coalescing to multi-agent simulations is \textbf{State Heterogeneity}. Micro-agents do not possess identical cognitive states: they maintain distinct psychological risk tolerances (Risk-Averse, Contrarian, Momentum), individualized cash and inventory balances ($\text{Inventory}_{i, t}$), and staggered execution arrival timestamps ($\Delta t_i \in [0.5\text{s}, 3.0\text{s}]$). If Singleflight attempted to hash the full agent decision prompt, every agent would produce a unique hash key, completely defeating request deduplication and unleashing a catastrophic GPU cache stampede.

To preserve the benefits of Singleflight while supporting full behavioral heterogeneity, AMPS strictly partitions the inference pipeline into two distinct phases:

\begin{enumerate}
    \item \textbf{Phase 1: Coalesced Macro News Parsing (Global Singleflight Pass):}
    When the 671B Macro Orchestrator broadcasts an unstructured regulatory filing or macroeconomic shock, all 400 micro-agents require an objective semantic understanding of the event. Rather than each agent independently querying the LLM to parse the headline, the first agent to process the shock acquires the Singleflight mutex. 
    
    The prompt consists strictly of the raw news text and objective market metrics (Order Flow Imbalance and market-wide volatility). The GPU executes a single forward pass, extracting an abstract, normalized 128-dimensional macroeconomic shock vector ($B_t \in \mathbb{R}^{128}$). When this single execution resolves, all waiting agent promises wake in $O(1)$ time, receiving the shared, zero-copy parsed shock vector.
    
    \item \textbf{Phase 2: Private Agent Policy Evaluation (Heterogeneous Vectorized Head):}
    Equipped with the shared 128D narrative shock vector $B_t$, each agent transitions to its private decision head. The agent concatenates $B_t$ with its private state variables -- local bid/ask queue position, net inventory holdings ($\text{Inventory}_{i, t}$), and psychological loss aversion thresholds -- into the 151-dimensional state vector $S_{i, t} \in \mathbb{R}^{151}$. 
    
    Because this stage evaluates lightweight 2-layer MLPs or quantized rules rather than dense autoregressive token generation, private evaluations execute concurrently across host CPU Gracemont E-cores without touching the GPU.
\end{enumerate}

\begin{figure}[htbp]
\centering
\begin{tikzpicture}[
    scale=0.88, transform shape,
    box/.style={draw, rectangle, rounded corners=3pt, align=center, font=\small, inner sep=5pt},
    phase1/.style={draw=blue!80!black, fill=blue!8, rectangle, rounded corners=4pt},
    phase2/.style={draw=green!70!black, fill=green!8, rectangle, rounded corners=4pt},
    arr/.style={->, thick, >=stealth}
]
    % Phase 1: Coalesced Macro News Parsing
    \node[box, fill=red!10, draw=red!70!black] (news) {Macroeconomic Shock\\ \footnotesize (SEC 8-K / Breaking Headline)};
    \node[box, fill=purple!15, draw=purple!80!black] (singleflight) [right=1.2cm of news] {Singleflight Mutex\\ \footnotesize 1st Caller Locks};
    \node[box, fill=orange!15, draw=orange!80!black] (gpu) [right=1.2cm of singleflight] {RTX 5070 Ti (vLLM)\\ \footnotesize 1x Macro Forward Pass};
    \node[box, fill=blue!15, draw=blue!80!black] (vector) [right=1.2cm of gpu] {128D Shock Vector\\ \footnotesize $B_t \in \mathbb{R}^{128}$ (Shared)};

    \draw[arr] (news) -- (singleflight);
    \draw[arr] (singleflight) -- node[above, font=\scriptsize] {1 Query} (gpu);
    \draw[arr] (gpu) -- node[above, font=\scriptsize] {Extract} (vector);

    % Phase 2: Private Agent Evaluation
    \node[box, fill=teal!10, draw=teal!70!black] (p_state) [below=1.6cm of singleflight] {Private Agent States\\ \footnotesize $\text{Inventory}_{i,t}$, Risk Limits, Queue};
    \node[box, fill=green!15, draw=green!80!black] (concat) [below=1.6cm of gpu] {Vector Assembly\\ \footnotesize $S_{i,t} = [L_t, \text{OFI}_t, S_t, \text{Inv}_{i,t}, B_t] \in \mathbb{R}^{151}$};
    \node[box, fill=yellow!20, draw=yellow!80!black] (ecores) [below=1.6cm of vector] {12 Gracemont E-Cores\\ \footnotesize INT8 VNNI Policy Passes};
    \node[box, fill=red!15, draw=red!80!black] (orders) [right=1.0cm of ecores] {Heterogeneous\\LOB Orders};

    \draw[arr, dashed, blue!80!black] (vector.south) -- ++(0,-0.6) -| (concat.north);
    \draw[arr] (p_state) -- (concat);
    \draw[arr] (concat) -- node[above, font=\scriptsize] {AVX2} (ecores);
    \draw[arr] (ecores) -- (orders);

    % Background groupings
    \begin{scope}[on background layer]
        \node[phase1, fit=(news)(singleflight)(gpu)(vector), label=above:{\textbf{\footnotesize Phase 1: Coalesced Macro News Parsing (Global GPU Pass)}}] {};
        \node[phase2, fit=(p_state)(concat)(ecores)(orders), label=below:{\textbf{\footnotesize Phase 2: Heterogeneous Private Agent Policy Evaluation (Concurrent CPU Passes)}}] {};
    \end{scope}
\end{tikzpicture}
\caption{Two-Phase Singleflight Architecture. Phase 1 coalesces macro shock comprehension into a single GPU forward pass, broadcasting an abstract 128D narrative vector $B_t$. Phase 2 evaluates heterogeneous agent actions concurrently across CPU Gracemont E-cores without touching GPU VRAM.}
\label{fig:two_phase_singleflight}
\end{figure}

This two-phase architecture eliminates over 95\% of redundant GPU token generation while fully preserving heterogeneous, individualized bounded rationality across the entire agent swarm.
